% Ссылка в тексте: (приложение~А)
\chapter{РУКОВОДСТВО ДЛЯ ИНТЕРВЬЮ С КЛИЕНТАМИ}

В~приложении представлено руководство для~проведения глубинных B2B-интервью
по~методу Mom~Test~\cite{fitzpatrick_mom_test_2013},
использовавшееся при~сборе первичных данных о~проблеме
(раздел~1.1 основного текста).

\section*{Контекст и цели интервью}

\textbf{Цель} — проверить гипотезу о~проблеме
«Разработка образовательных курсов занимает слишком много времени и/или
требует компетенций, которых нет у~эксперта предметной области».

\textbf{Целевая аудитория} — HR-специалисты, L\&D-менеджеры,
методисты, тренеры в~B2B-компаниях с~>~50~сотрудников.

\textbf{Формат} — 30–45~минут, видеозвонок или~очно.

\section*{Вопросы руководства}

\textbf{Блок 1. Разогрев и контекст}
\begin{enumerate}
  \item Расскажите о~своей роли. С какой частотой вам приходится создавать
    учебные материалы для~сотрудников?
  \item Какой последний курс или~материал вы разрабатывали?
    Вспомните этот процесс — с~чего он начался?
  \item Сколько времени ушло от~идеи до~готового материала?
\end{enumerate}

\textbf{Блок 2. Проблема}
\begin{enumerate}
  \setcounter{enumi}{3}
  \item Что в~этом процессе было самым сложным или~утомительным?
  \item Были ли случаи, когда вы \textit{не~смогли} создать нужный
    материал из-за нехватки времени, ресурсов или~экспертизы?
  \item С какой частотой учебные материалы устаревают и~требуют обновления?
    Каким образом вы сейчас с~этим справляетесь?
\end{enumerate}

\textbf{Блок 3. Текущие решения}
\begin{enumerate}
  \setcounter{enumi}{6}
  \item Какие инструменты или~сервисы вы ранее пробовали для~автоматизации
    разработки курсов? Что не~устроило?
  \item Если бы вы могли потратить бюджет на~решение этой проблемы~—
    на~что бы потратили?
\end{enumerate}

\textbf{Блок 4. Закрытие}
\begin{enumerate}
  \setcounter{enumi}{8}
  \item Есть ли коллеги, которых стоит поговорить на~эту тему?
  \item Можно ли мне прислать вам краткий обзор после дополнительных
    разговоров с~ещё несколькими людьми?
\end{enumerate}

\section*{Сводка ключевых выводов}

По~результатам 20 проведённых интервью устойчивее всего повторялись
следующие сигналы.

\begin{table}[H]
  \caption{Агрегированные выводы интервью с клиентами}
  \label{tab:custdev}
  \centering
  \begin{tabular}{|P{5cm}|P{2.5cm}|P{5.5cm}|}
  \hline
  \textbf{Наблюдение} & \textbf{Значение} & \textbf{Интерпретация} \\
  \hline
  Время разработки — главное ограничение
    & 70\,\%
    & проблема носит системный характер \\
  \hline
  Используют подрядчиков / внешних методистов
    & 50\,\%
    & рынок готов платить за~снижение этой нагрузки \\
  \hline
  Готовы рассмотреть ИИ-инструмент
    & 60\,\%
    & зафиксирован ранний интерес к~решению \\
  \hline
  Формат исследования
    & 20 интервью
    & глубинные B2B-интервью с клиентами по~методу Mom Test \\
  \hline
  \end{tabular}
\end{table}

\begin{itemize}
  \item Наибольшую ценность респонденты видят в~ускорении подготовки
    первого черновика курса.
  \item Для~нормативных и~регламентных сценариев критична прослеживаемость
    к~источнику и~возможность ручной проверки.
  \item Для~B2B-сегмента значимы суверенитет данных и~возможность
    локального развёртывания.
\end{itemize}
